%%%%%%%%%%%%%%%%%%%%%%%%%%%%%%%%%%%%%%%%%%%%%%%%%%%%%%%%%%%%%%%%%%%%%%%%%%%%%%
% FireSpreadNet – A0 Landscape Poster (single page)
%%%%%%%%%%%%%%%%%%%%%%%%%%%%%%%%%%%%%%%%%%%%%%%%%%%%%%%%%%%%%%%%%%%%%%%%%%%%%%

\documentclass[a0,landscape]{a0poster}

%----------------------------------------------------------------------------------------
% PACKAGES
%----------------------------------------------------------------------------------------
\usepackage[svgnames]{xcolor}
\usepackage{multicol}
\columnsep=60pt
\columnseprule=1.5pt
\usepackage[absolute,overlay]{textpos}
\usepackage{times}
\usepackage{graphicx}
\graphicspath{{figures/}{Images/}}
\usepackage{booktabs}
\usepackage[font=small,labelfont=bf]{caption}
\usepackage{amsfonts, amsmath, amsthm, amssymb}
\usepackage{url}
\usepackage{tcolorbox}
\tcbuselibrary{skins,breakable}
\usepackage{array}
\usepackage{tabularx}
\usepackage{enumitem}

%----------------------------------------------------------------------------------------
% COLOR PALETTE
%----------------------------------------------------------------------------------------
\definecolor{FireRed}{RGB}{180, 40, 20}
\definecolor{FireOrange}{RGB}{220, 100, 20}
\definecolor{SmokeGrey}{RGB}{90, 90, 90}
\definecolor{AshLight}{RGB}{245, 240, 235}
\definecolor{ForestGreen}{RGB}{34, 100, 45}
\definecolor{DeepBlue}{RGB}{25, 60, 120}
\definecolor{AlertYellow}{RGB}{255, 200, 0}
\definecolor{LightGrey}{RGB}{230, 230, 230}

%----------------------------------------------------------------------------------------
% CUSTOM COMMANDS
%----------------------------------------------------------------------------------------
\newcommand{\sectitle}[1]{%
  \par\noindent\colorbox{FireRed!15}{%
    \parbox{\linewidth}{%
      \vspace{0.12em}%
      {\Large\bfseries\color{FireRed} \hspace{0.3em}#1}%
      \vspace{0.12em}%
    }%
  }\par\vspace{0.1em}%
}

\newcommand{\subsectitle}[1]{%
  \par\noindent{\large\bfseries\color{DeepBlue} #1}\par\vspace{0.1em}%
}

\newcommand{\keybox}[1]{%
  \begin{tcolorbox}[
    colback=FireRed!10,colframe=FireRed,
    arc=3pt,boxrule=1pt,
    left=4pt,right=4pt,top=3pt,bottom=3pt
  ]{\small #1}\end{tcolorbox}%
}

\newcommand{\greenbox}[1]{%
  \begin{tcolorbox}[
    colback=ForestGreen!10,colframe=ForestGreen,
    arc=3pt,boxrule=1pt,
    left=4pt,right=4pt,top=3pt,bottom=3pt
  ]{\small #1}\end{tcolorbox}%
}

\newcommand{\bluebox}[1]{%
  \begin{tcolorbox}[
    colback=DeepBlue!8,colframe=DeepBlue,
    arc=3pt,boxrule=1pt,
    left=4pt,right=4pt,top=3pt,bottom=3pt
  ]{\small #1}\end{tcolorbox}%
}

\begin{document}
\enlargethispage{200pt}

%========================================================================================
% HEADER
%========================================================================================
\begin{minipage}[b]{0.76\linewidth}
  \begin{flushleft}
    {\fontsize{72}{80}\selectfont\bfseries\color{FireRed} FireSpreadNet}\\[0.15cm]
    {\LARGE\color{DeepBlue}\textit{Attention-Enhanced Deep Learning for Next-Day Wildfire Spread Prediction}}\\[0.12cm]
    {\large\color{SmokeGrey}
      Benchmark Study with Physics-Informed Hypothesis, Explainability (XAI), and Deployment}\\[0.15cm]
    {\normalsize\color{SmokeGrey}
      \textbf{Syrine Maaref \& Dorra Trabelsi}\quad|\quad
      Supervisor: Ahmed Ben Taleb\quad|\quad ITBS}
  \end{flushleft}
\end{minipage}
%
\begin{textblock*}{12cm}(93cm,4cm)
  \includegraphics[width=22cm]{logoitbs.png}
\end{textblock*}

\vspace{0.4cm}

%========================================================================================
% TOP FULL-WIDTH BLOCK: Motivation + Problem Statement
%========================================================================================
\noindent
\begin{tcolorbox}[
  enhanced,
  colback=AshLight,
  colframe=FireRed,
  arc=5pt,
  boxrule=2pt,
  left=10pt, right=10pt, top=6pt, bottom=6pt,
  title={\Large\bfseries\color{white} Motivation \& Problem Statement},
  attach boxed title to top left={yshift=-2mm, xshift=8mm},
  boxed title style={colback=FireRed, colframe=FireRed, arc=3pt}
]
\begin{minipage}[t]{0.54\linewidth}
  \vspace{0pt}
  {\small\color{SmokeGrey}%
    Wildfires burned more than \textbf{420 million hectares} in 2023--2024, causing massive human and
    ecological losses. Emergency response remains largely \textbf{reactive}: resources are deployed
    \emph{after} ignition, limiting evacuation time and increasing risk.
    Satellite imagery now provides daily multimodal observations at scale,
    enabling data-driven models that are faster and more adaptive than traditional physical simulators,
    yet still explainable enough to support expert decision-making.%
  }
\end{minipage}
\hfill\vrule width 1.5pt\hfill
\begin{minipage}[t]{0.42\linewidth}
  \vspace{0pt}
  \begin{tcolorbox}[
    colback=DeepBlue!12, colframe=DeepBlue,
    arc=4pt, boxrule=2pt,
    left=6pt, right=6pt, top=5pt, bottom=5pt
  ]
  \begin{center}
    {\large\bfseries\color{DeepBlue} Research Question}\\[0.3em]
    {\normalsize\color{black}%
      \textit{Can attention-based deep learning on multimodal satellite data}\\
      \textit{reliably predict next-day wildfire spread, outperform physical baselines,}\\
      \textit{and remain explainable enough for operational expert use?}%
    }
  \end{center}
  \end{tcolorbox}
\end{minipage}
\end{tcolorbox}

\vspace{0.35cm}

%========================================================================================
% THREE-COLUMN BODY
%========================================================================================
\begin{multicols}{3}
\raggedcolumns
\color{SmokeGrey}
\small

%----------------------------------------------------------------------
% COLUMN 1 – Dataset
%----------------------------------------------------------------------
\sectitle{Dataset: Next-Day Wildfire Spread}

\begin{itemize}[leftmargin=1.2em,itemsep=0.05em]
  \item \textbf{18,545 patches}, 64$\times$64 px ($\approx$1 km/px) — continental USA.
  \item \textbf{12 input channels}: fire history, weather (wind, humidity, temp., rain), drought (PDSI),
        vegetation (NDVI, ERC), topography (elevation, slope), population density.
  \item \textbf{Target}: binary next-day fire mask.
  \item \textbf{Temporal split} (no leakage): 14,979 train / 1,877 val / 1,689 test.
  \item \textbf{Class imbalance}: only \textbf{1.07\%} fire pixels.
\end{itemize}

\vspace{0.15em}
\keybox{%
  \textbf{Imbalance mitigation:}
  Focal Loss ($\alpha=0.989$, derived analytically) at pixel level
  + \texttt{WeightedRandomSampler} at sample level to over-represent fire patches.
}

\vspace{0.2em}
\begin{center}
  \includegraphics[width=0.88\linewidth]{Images/eda_sample_9804.png}\\[0.15em]
  {\footnotesize\textbf{Figure 1.} Example 12-channel wildfire patch.}
\end{center}

\vspace{0.2em}
\begin{center}
  \includegraphics[width=0.88\linewidth]{Images/eda_split_comparison.png}\\[0.15em]
  {\footnotesize\textbf{Figure 2.} Train / val / test split distribution.}
\end{center}

\columnbreak

%----------------------------------------------------------------------
% COLUMN 2 – Architecture
%----------------------------------------------------------------------
\sectitle{Architecture: U-Net + Attention (Best Model)}

\vspace{0.1em}
\bluebox{%
  Four models were benchmarked under \textbf{identical conditions} (same split, optimizer AdamW,
  Focal+Dice loss, seed). Only U-Net+Attention is detailed here as the top performer.
}

\vspace{0.2em}
\subsectitle{U-Net with Attention Gates}

The encoder--decoder backbone follows the U-Net design. \textbf{Attention gates} are inserted at each skip connection to suppress irrelevant spatial activations and focus the decoder on fire-prone regions:

$$\boldsymbol{\alpha}^l \;=\;
  \sigma\!\Bigl(\mathbf{W}_\psi\,
    \mathrm{ReLU}\!\bigl(\mathbf{W}_x\mathbf{x}^l + \mathbf{W}_g\mathbf{g}^l + b\bigr)
  + b_\psi\Bigr)$$

where $\mathbf{x}^l$ is the skip-connection feature, $\mathbf{g}^l$ is the gating signal from the
lower level, and $\boldsymbol{\alpha}^l\!\in(0,1)$ is the soft spatial mask applied before
concatenation.

\vspace{0.15em}
\begin{center}
  % FIGURE RECOMMENDATION: U-Net+Attention diagram (encoder, bottleneck, decoder, attention gates)
  \includegraphics[width=0.88\linewidth]{Images/training_curves_local.png}\\[0.15em]
  {\footnotesize\textbf{Figure 3.} Training / validation loss and Dice convergence.}
\end{center}

\vspace{0.15em}
\begin{center}
{\footnotesize
\renewcommand{\arraystretch}{1.2}
\begin{tabular}{lcc}
\toprule
\textbf{Component} & \textbf{Detail} \\
\midrule
Backbone & U-Net (4 enc.\ levels) \\
Skip connections & Attention-gated \\
Parameters & $\sim$1.97M \\
Loss & 0.5 Focal + 0.5 Dice \\
Optimizer & AdamW, lr = 1e-4 \\
Threshold $\tau^*$ & 0.90 (val-tuned) \\
\bottomrule
\end{tabular}
}\end{center}

\columnbreak

%----------------------------------------------------------------------
% COLUMN 3 – Results + XAI + Conclusion
%----------------------------------------------------------------------
\sectitle{Results}

\begin{center}
{\footnotesize
\renewcommand{\arraystretch}{1.2}
\begin{tabular}{lcccc}
\toprule
\textbf{Model} & \textbf{Dice}$\uparrow$ & \textbf{IoU}$\uparrow$ & \textbf{Recall}$\uparrow$ & \textbf{AUC-ROC}$\uparrow$ \\
\midrule
CA (physics)          & 0.3254 & 0.1943 & 0.4734 & -- \\
ConvLSTM              & 0.3185 & 0.1894 & 0.4865 & -- \\
PI-CCA                & 0.3240 & 0.1933 & 0.3767 & -- \\
\textbf{U-Net+Attn}   & \textbf{0.4034} & \textbf{0.2526} & \textbf{0.5418} & \textbf{0.9415} \\
\midrule
Huot et al.\ (2022)   & 0.311  & --     & --     & 0.839 \\
\bottomrule
\end{tabular}
}
\end{center}

\vspace{0.05em}
\greenbox{%
  \textbf{H1 confirmed:} U-Net+Attention achieves \textbf{Dice = 0.4034}, \textbf{IoU = 0.2526},
  and \textbf{AUC-ROC = 0.9415} — outperforming all baselines and the Huot et al.\ (2022)
  published benchmark.
}
\vspace{0.05em}
\keybox{%
  \textbf{Physics prior not validated:} PI-CCA shows lower recall than U-Net+Attention and
  does not outperform it, suggesting the CA prior introduces more bias than benefit at this scale.
}

\vspace{0.15em}
\sectitle{Explainability (XAI)}

\begin{center}
  \includegraphics[width=0.72\linewidth]{Images/shap_importance_unet.png}\\[0.1em]
  {\footnotesize\textbf{Figure 4.} SHAP global feature importance (U-Net+Attention).}
\end{center}

\vspace{0.05em}
\bluebox{%
  \textbf{prev\_fire\_mask} (fire history) dominates SHAP importance ($\approx$80\%),
  consistent with the physical persistence of fire perimeters.
  Wind and drought indices rank next, confirming meteorological drivers.
  Grad-CAM spatial maps reveal that attention gates concentrate on active-fire boundaries.
}

\vspace{0.15em}
\sectitle{Deployment \& Conclusion}

\textbf{Deployment:} Dockerized inference API + expert-facing web dashboard.
Safety guardrails: probability output only, threshold under expert control, no autonomous alerts.

\vspace{0.1em}
\greenbox{%
  \textbf{Conclusion:} U-Net+Attention with spatial attention gates reaches
  \textbf{Dice = 0.4034 / IoU = 0.2526 / AUC-ROC = 0.9415},
  beating physical simulators and prior deep-learning baselines.
  SHAP \& Grad-CAM make predictions explainable.
  \textit{Future work}: stronger physics priors, uncertainty quantification, multi-step horizon.
}

\end{multicols}

\end{document}
